\chapter{把整条证明重新连起来}\label{sec:completion}
\section{现在可以不用公式，先把逻辑讲一遍}
起点是任意有理配置。先减去均值、清分母、约去公因子，得到本原整数零和代表。

如果 $G$ 含非 $p$ 素因子，那个素数下的共同非零余数永远不变，立即判不可解。若 $G$ 只有 $p$ 因子，我们按位置数选择证明路线。

位置足够多，保持两个重复块，选择安全整数平均，使整数平方和下降；旧容量不等式遗漏的窄段，则利用其大量素因子带来的强计数估计另行补齐。

临界带位置紧张，先安全地压缩成三个块。通过真实返回得到剪切、倍率和逆，进而取得足够大的主同余子群。有限层面将参数与某个零和终端列对齐，再由主同余子群把它们精确对齐。最后真的平均一个零和组，接含零收尾。

小素数 $3,5,7$ 中一般容量不适用的部分，由前两讲的独立下降或有限证书补上。每一块都提供指定范围内的全部输入结论，下面把范围合并。
\begin{proof}[定理\ref{thm:prime}的证明]
必要性由引理\ref{lem:obstruction}给出，下界由命题\ref{prop:lower}给出。对于充分性，$p\ge11$ 时：$n\ge4p$ 用命题\ref{prop:wide}；$n=3p$ 用命题\ref{prop:threep}；其余临界带先用命题\ref{prop:entry}压缩，再按下表完成。
\begin{center}
\begin{tabular}{ll}
\toprule 参数范围 & 核内结论\\\midrule
$n=2p+r,\ 1\le r<p,\ r$ 奇 & 命题\ref{prop:odd}\\
$r=2$ & 命题\ref{prop:r2}\\
$4\le r<p,\ r$ 偶 & 命题\ref{prop:even}\\
$3p<n<4p$ & 命题\ref{prop:upper}\\\bottomrule
\end{tabular}
\end{center}
$p=5$ 时，$11\le n\le15$ 已由上述下半带及等块结论处理；$n=16,17,18$ 用引理\ref{lem:collision-small}；$n=19$ 满足引理\ref{lem:integer-range}；$n\ge20$ 用命题\ref{prop:wide}。
$p=7$ 时，下半带处理 $15\le n\le21$；$n=22$ 用引理\ref{lem:seven22}；$23\le n\le26$ 用引理\ref{lem:collision-small}；$n=27$ 满足引理\ref{lem:integer-range}；更大维数用命题\ref{prop:wide}。
$p=3$ 时，$n=7$ 用命题\ref{prop:entry}及\ref{prop:odd}，$n=8,10$ 用引理\ref{lem:eight}、\ref{lem:ten}，$n\ge11$ 用引理\ref{lem:triple-large}。$n=9$ 将位置排成 $3\times3$，先平均各行再平均各列，六次操作即达到全局均值。
这些范围覆盖全部奇素数及全部 $n\ge2p+1$。推论\ref{cor:threshold}随定理和第\ref{sec:obstruction}讲的二平均论证成立。
\end{proof}

\section{如果有人问：你们到底用了多深的数学}
必要条件用最大公因数和模运算；整数路线用重复值、零和选择、抽屉原理和平方和；压缩用线性组合和CRT；核心矩阵的制造用真实整数计数、标量周期和交换子。

最深的外部输入出现在从“已经实现上下两整族剪切”到“已经实现某个主同余子群”的那一步，使用 Morris 和 Serre 的定理。取得它以后，把同余相同的两列精确连接，又是裴蜀和二阶矩阵的直接计算。

所以整个证明既不是纯粹平均值不等式，也不是套一个数论黑箱就结束。具体的位置构造提供可用矩阵，数论把这些矩阵的能力扩展到任意有理参数，最后仍要回到实际位置归零。

\section{读完以后，最值得自己回答的几个问题}
\begin{enumerate}
\item 为什么 $p$ 和 $n$ 的素因子承担不同角色？答：$p$ 决定实际分母，非 $p$ 的 $n$ 素因子承担共同余数障碍。
\item 为什么平方和下降不总能证明有限终止？答：有理正数可以无限下降；整数格才给离散性。
\item 为什么可逆矩阵不直接把列送到零？答：它只先送到含零和组的非零核心，最后由奇异平均归零。
\item 为什么同余子群出现以后仍需要有限计算？答：它能在一个剩余类里精确输送，但初始剩余类还未必是终端类。
\item 为什么证明还分很多计数情形？答：不同块重数提供的真实副本数量不同；同一抽象矩阵不一定在每种库存下都能实现。
\end{enumerate}
这些问题说明了判据简短而充分性很长的原因。公式 $G=p^e$ 只描述障碍；要证明它也是充分条件，还必须造出不会丢掉余数见证、能在有限步内到达终点的路径。
